\chapter{System design}
\label{cha:design}

This chapter presents the architecture that follows from the requirements in \Cref{cha:vision} and the technology choices in \Cref{cha:methods-and-tools}, leaving the difficulties met while building it to \Cref{cha:implementation}. Its organising principle is a division of labour by trustworthiness: the language model will be given the work that requires interpretation -- reading a request, choosing dishes, writing a recipe, explaining a rationale -- and will be denied the work where a plausible-looking wrong answer would be indistinguishable from a right one. Two structural devices will enforce that boundary: an output schema with no field in which a nutrient value could be written, and a resolution step that turns the model's food references into knowledge base records before any arithmetic happens. The knowledge base, retrieval strategy, totaller and review loop will follow from requiring those two devices to hold in every configuration the experiments switch between.

\section{Domain model: references versus hydrated plans}
\label{sec:design-domain-model}

A meal plan will take two distinct representations inside the system, and the boundary between them is where the no-invented-foods contract will be enforced.

\begin{definition}[Agent plan]
    \label{def:agent-plan}
    An \emph{agent plan} is the structured output a language model produces: a set of meals, each a named recipe with a list of portions, where every portion is a triple of a database identifier, a human-readable name and a quantity in grams. It carries no nutrient field of any kind.
\end{definition}

\begin{definition}[Hydration and hydrated plans]
    \label{def:hydrated-plan}
    \emph{Hydration} resolves every identifier in an agent plan against the real food knowledge base and replaces it with that food's full database record, carrying its actual nutrient profile; the term is used in the sense it carries in data-mapping nomenclature, where hydrating a reference means filling it in with the record it points to. A \emph{hydrated plan} is the result of applying hydration to an agent plan.
\end{definition}

\begin{figure}[p]
    \centering
    \begin{tikzpicture}[
    x=1mm, y=1mm,
    proc/.style={draw, rounded corners, fill=white, anchor=center, text width=42mm,
                 minimum height=12mm, inner sep=1.8mm, font=\small, align=center},
    data/.style={draw, dashed, fill=white, align=left,
                 text width=52mm, inner sep=2.2mm, font=\small},
    entry/.style={draw, dashed, fill=white, align=left,
                  text width=38mm, inner sep=2mm, font=\small},
    store/.style={draw, rounded corners, fill=white, anchor=center, text width=24mm,
                  minimum height=28mm, inner sep=1.8mm, font=\small, align=center},
    arr/.style={-{Latex[length=2mm]}},
    dcall/.style={-{Latex[length=2mm]}, draw=black!55, dashed},
    edgelab/.style={font=\footnotesize, text=black!60, align=center},
    lin/.style={},
]

\node[proc] (nutrition) at (32, 0) {Nutrition agent};

\node[proc] (tool) at (104, 0) {Totaller tool};

\node[data] (plan) at (32, -26) {
    \textbf{Agent plan}\\
    rationale,\\
    meals: list of portions + recipe
};

\node[proc] (hydration) at (104, -26) {Hydration};

\node[data] (hplan) at (104, -54) {
    \textbf{Hydrated plan}\\
    rationale,\\ 
    meals: list of portions with nutrients per-\qty{100}{\gram} + recipe
};

\node[proc] (totaller) at (104, -88) {Totaller};

\node[data] (totals) at (104, -112) {
    \textbf{Nutrient totals}\\
    day and per-meal totals,\\
    coverage caveats
};

\node[store] (kb) at (32, -64) {Food\\knowledge\\base};

\node[entry] (dbentry) at (32, -94) {
    \textbf{Base entry}\\
    identifier,\\
    name,\\
    nutrients per-\qty{100}{\gram}
};

\draw[arr] (32, 16) -- (nutrition.north);
\draw[arr] (totals.south) -- ++(0, -12);

\draw[lin] (nutrition) -- (plan);
\draw[arr] (plan.east) -- (hydration.west);
\draw[lin] (hydration) -- (hplan);
\draw[arr] (hplan.south) -- (totaller.north);
\draw[lin] (totaller) -- (totals);

\draw[dcall] ([yshift=3.5mm]nutrition.east) -- node[edgelab, above] {agent plan} ([yshift=3.5mm]tool.west);
\draw[dcall] ([yshift=-3.5mm]tool.west) -- node[edgelab, below] {totals} ([yshift=-3.5mm]nutrition.east);
\draw[arr] (tool.south) -- (hydration.north);
\draw[dcall] (totals.east) -- (146, -112) -- (146, 0) -- (tool.east);

\draw[lin] (kb.south) -- (dbentry.north);
\draw[arr] (dbentry.east) -- (68, -94) -- (68, -30) -- ([yshift=-4mm]hydration.west);

\end{tikzpicture}

    \caption{Resolution of portion references against the food knowledge base, and the route the agent-facing totaller tool takes through it. Dashed edges are the tool round trip.}
    \label{fig:diagram-hydration}
\end{figure}

The totaller will always operate on a hydrated plan (\Cref{def:hydrated-plan}). What the model hands it mid-generation, however, will still be an agent plan (\Cref{def:agent-plan}), so the two will be kept apart, as \Cref{fig:diagram-hydration} shows: the tool the agent can call is a thin wrapper that first pushes the draft through the same hydration step the rest of the pipeline uses and only then passes the result to the totaller. The nutrient figures the model gets back will therefore be computed from database records and not from anything it claims. This separation is what will make the no-invented-foods contract in \Cref{sec:design-no-invented-foods} structurally enforceable instead of merely a matter of prompt wording: an agent plan will have no field in which a nutrient value could even be written down, so there will be nothing for the model to hallucinate.

\section{The no-invented-foods contract}
\label{sec:design-no-invented-foods}

\ref{req:no-invented-nutrients}~requires that the system never report a nutrient figure that cannot be traced to a database entry and a specific quantity. This will be enforced at two independent points. First, the agent plan's schema (\Cref{sec:design-domain-model}) will simply have no place to put a nutrient value. A portion will be a code, a name and a gram quantity, so a model wishing to report an invented calorie count will only be able to write it into free text, which will never be parsed as nutrition data. Second, every generated plan will be validated against the real food knowledge base before it is accepted, attempting to resolve every referenced code. A validator failure will raise a structured retry request back to the model, listing exactly which codes did not resolve and reminding it that every code must be copied verbatim from a prior lookup result instead of guessed. Every accepted plan will therefore pass this code check, so the no-invented-foods guarantee will hold identically regardless of which code path produced the plan.

\section{Architecture overview}
\label{sec:design-pipeline}

The system will be structured as a single, mostly linear pipeline invoked once per session, not as a long-running stateful service; \Cref{fig:diagram-pipeline} shows its stages. A runner will start the session, fetching the profile and its frozen targets, reading the variant flags that decide which optional mechanisms are active (\Cref{sec:design-variant-config}), and resolving a single model specification shared by the nutrition agent, the critic and the refiner. The nutrition agent will compose the day's recipes and portions, drawing on the food knowledge base (\Cref{sec:design-food-db}) and, when it is enabled, the totaller (\Cref{sec:design-totaller}). Its output will be an agent plan (\Cref{def:agent-plan}), which the output validator (\Cref{sec:design-no-invented-foods}) will check against the real database; when the critique-and-refine loop is enabled the validated plan will pass through review (\Cref{sec:design-reflection}) before hydration (\Cref{def:hydrated-plan}) turns it into the delivered meal plan.

\begin{figure}[p]
    \centering
    \resizebox{\textwidth}{!}{\begin{tikzpicture}[
    x=1mm, y=1mm,
    box/.style={draw, rounded corners, anchor=center, text width=34mm,
                minimum height=11mm, inner sep=1.2mm, font=\scriptsize, align=center},
    agent/.style={box, minimum height=20mm},
    side/.style={draw, rounded corners, anchor=center, text width=22mm,
                 minimum height=9mm, inner sep=1.2mm, font=\scriptsize, align=center},
    lane/.style={draw, rounded corners, anchor=center, text width=18mm,
                 inner sep=1.2mm, font=\scriptsize, align=center},
    dcall/.style={-{Latex[length=1.6mm]}, draw=black!55, dashed},
    ret/.style={-{Latex[length=1.6mm]}, draw=black!55, dashed},
    ensured/.style={{Latex[length=1.6mm]}-{Latex[length=1.6mm]}},
    edgelab/.style={font=\scriptsize, text=black!60, align=center},
    arr/.style={-{Latex[length=2mm]}},
    % validator marker: a hollow triangle sitting on an edge, deliberately
    % unlike the solid arrow heads used for flow
    pics/vtri/.style={code={\filldraw[fill=white, draw=black]
        (-1.8, -2.2) -- (-1.8, 2.2) -- (2.0, 0) -- cycle;}},
]

% --- centre column (x = 48) ---
\node[box] (runner) at (48, 0) {Runner};
\node[agent] (nutrition) at (48, -30) {Nutrition agent};
\node[agent] (critic) at (48, -70) {Critic};
\node[agent] (refiner) at (48, -110) {Refiner};
\node[agent] (hydration) at (48, -144) {Hydration};

% --- profiles store, top left ---
\node[side] (profiles) at (0, 0) {Profiles store};

% --- tool lanes on opposite sides ---
\node[lane, minimum height=140mm] (lookup) at (-10, -85) {Food knowledge database};
\node[lane, minimum height=108mm] (totaller) at (106, -70) {Totaller};

% --- attachment points, taken from the block borders themselves ---
\coordinate (n_call_w) at ([yshift=6mm]nutrition.west);
\coordinate (n_ret_w)  at ([yshift=0mm]nutrition.west);
\coordinate (n_retry)  at ([yshift=-6mm]nutrition.west);
\coordinate (n_call_e) at ([yshift=6mm]nutrition.east);
\coordinate (n_ret_e)  at ([yshift=0mm]nutrition.east);
\coordinate (r_call_w) at ([yshift=4mm]refiner.west);
\coordinate (r_ret_w)  at ([yshift=-4mm]refiner.west);
\coordinate (r_call_e) at ([yshift=4mm]refiner.east);
\coordinate (r_ret_e)  at ([yshift=-4mm]refiner.east);
\coordinate (c_out)    at ([yshift=-6mm]critic.west);
\coordinate (h_in)     at ([yshift=6mm]hydration.west);

% --- main vertical flow ---
\draw[arr] (48, 14) -- node[edgelab, right] {prompt} (runner);
\draw[arr] (runner) -- node[edgelab, right] {runs} (nutrition);
\pic[rotate=-90] at ([yshift=-2mm]nutrition.south) {vtri};
\draw[arr] ([yshift=-4mm]nutrition.south) -- node[edgelab, right] {plan} (critic.north);
\draw[arr] (hydration) -- node[edgelab, right] {hydrated plan} (48, -160);

% --- critic and refiner: two parallel edges, opposite directions ---
\pic[rotate=-90] at ([xshift=-6mm, yshift=-2mm]critic.south) {vtri};
\draw[arr] ([xshift=-6mm, yshift=-4mm]critic.south) -- node[edgelab, anchor=east, pos=0.55] {feedback} ([xshift=-6mm]refiner.north);
\pic[rotate=90] at ([xshift=6mm, yshift=2mm]refiner.north) {vtri};
\draw[arr] ([xshift=6mm, yshift=4mm]refiner.north) -- node[edgelab, anchor=west, pos=0.55] {refined plan} ([xshift=6mm]critic.south);

% --- approved plan: critic -> hydration ---
\draw[arr] (c_out) -- (14, -76) -- (14, -138) -- (h_in);
\node[edgelab, anchor=south] at (20, -75) {approved plan};

% --- profiles feed ---
\draw[arr] (profiles) -- node[edgelab, above] {profile} (runner);

% --- left lane: food knowledge database (call validated on the way in) ---
\pic[rotate=180] at ([xshift=-2mm]n_call_w) {vtri};
\draw[dcall] ([xshift=-4mm]n_call_w) -- node[edgelab, above, pos=0.7] {tool call} (n_call_w -| lookup.east);
\draw[ret] (n_ret_w -| lookup.east) -- node[edgelab, below, pos=0.3] {tool response} (n_ret_w);
\pic[rotate=180] at ([xshift=-2mm]r_call_w) {vtri};
\draw[dcall] ([xshift=-4mm]r_call_w) -- node[edgelab, above, pos=0.7] {tool call} (r_call_w -| lookup.east);
\draw[ret] (r_ret_w -| lookup.east) -- node[edgelab, below, pos=0.3] {tool response} (r_ret_w);

% hydration <-> food knowledge database: deterministic, no validator
\coordinate (h_dbc) at ([yshift=0mm]hydration.west);
\coordinate (h_dbr) at ([yshift=-6mm]hydration.west);
\draw[arr] (h_dbc) -- node[edgelab, above, pos=0.7] {call} (h_dbc -| lookup.east);
\draw[-{Latex[length=1.6mm]}] (h_dbr -| lookup.east) -- node[edgelab, below, pos=0.3] {result} (h_dbr);

% --- right lane: totaller (call validated on the way in) ---
\pic at ([xshift=2mm]n_call_e) {vtri};
\draw[dcall] ([xshift=4mm]n_call_e) -- node[edgelab, above, pos=0.5] {tool call} (n_call_e -| totaller.west);
\draw[ret] (n_ret_e -| totaller.west) -- node[edgelab, below, pos=0.5] {tool response} (n_ret_e);
\pic at ([xshift=2mm]r_call_e) {vtri};
\draw[dcall] ([xshift=4mm]r_call_e) -- node[edgelab, above, pos=0.5] {tool call} (r_call_e -| totaller.west);
\draw[ret] (r_ret_e -| totaller.west) -- node[edgelab, below, pos=0.5] {tool response} (r_ret_e);

% critic <-> totaller: call validated on the way in
\coordinate (c_call_e) at ([yshift=4mm]critic.east);
\coordinate (c_ret_e)  at ([yshift=-4mm]critic.east);
\draw[arr] (c_call_e) -- node[edgelab, above, pos=0.5] {call} (c_call_e -| totaller.west);
\draw[-{Latex[length=1.6mm]}] (c_ret_e -| totaller.west) -- node[edgelab, below, pos=0.5] {result} (c_ret_e);

% --- legend ---
\pic at (-19, -166) {vtri};
\node[edgelab, anchor=west, text=black] at (-15, -166) {pydantic-ai validator};

\end{tikzpicture}
}
    \caption{End-to-end pipeline in its full configuration.}
    \label{fig:diagram-pipeline}
\end{figure}

\section{Food knowledge base}
\label{sec:design-food-db}

The food knowledge base will combine the two bulk data sources chosen in \Cref{sec:tools-food-data} behind a single lookup interface, stored in DuckDB (\Cref{sec:tools-storage}). Each source will be materialised into its own flat table, since their native schemas differ too much to usefully merge: a packaged product's brand and ingredient list has no equivalent in a generic USDA entry, and vice versa. A fixed set of nutrients will be carried for every entry: energy, the three macronutrients (protein, carbohydrate and fat), and eleven further fields -- saturated fat, fibre, sugar, salt, sodium, potassium, calcium, iron, vitamin~C, vitamin~D and cholesterol. These are the nutrients both sources report often enough to be worth tracking, and they are what every total in this thesis is drawn from. Every nutrient column will be normalised to a common unit and expressed per~\qty{100}{\gram} of food, so that scaling by an arbitrary portion size is a single multiplication regardless of source. Every entry's identifier will be namespaced by its source -- an Open Food Facts barcode or a USDA FoodData Central identifier -- so that the two tables can share a single reference format without collision. Each table will carry a full-text search index over its name and description fields and an embedding vector per entry, computed at build time with the embedding model chosen in \Cref{sec:tools-embeddings}, to support the hybrid search described next.

\section{Hybrid search and the lookup tool contract}
\label{sec:design-hybrid-search}

The lookup tool exposed to the nutrition agent will accept a batch of natural-language queries in a single call. Each query will carry a separate cap on how many hits it wants from each data source, and by default it will be run against both, and within each source against both the lexical and the dense retrieval channels chosen in \Cref{sec:tools-retrieval}, fused by reciprocal rank fusion into one ranked list per source, returned side by side and not merged across sources (\Cref{fig:diagram-retrieval}). Setting a source's cap to zero will skip it, so a query whose intent is clearly branded or clearly generic can be aimed at one catalogue; the prompt will ask for that where it applies. Keeping the two catalogues unmerged is deliberate: a branded Open Food Facts hit and a generic USDA hit for the same query are not interchangeable, and a single fused ranking would hide which source a code came from. Each hit will carry its identifier, name and per-\qty{100}{\gram} nutrient profile, so the model can compare candidates and choose a portion size without a further round trip. Why the tool is batched, how many hits it returns, and what a single call costs against the context budget are implementation decisions, taken up in \Cref{sec:impl-tool-ergonomics}.

\begin{figure}[htbp]
    \centering
    \begin{tikzpicture}[
    node distance=7mm and 6mm,
    box/.style={draw, rounded corners, align=center, text width=26mm, minimum height=9mm, inner sep=1.5mm, font=\small},
    band/.style={box, text width=61mm},
    call/.style={draw, rounded corners, align=center, text width=50mm, minimum height=9mm, inner sep=1.5mm, font=\small},
    resultbox/.style={draw, rounded corners, align=center, text width=56mm, minimum height=9mm, inner sep=1.5mm, font=\small},
    frame/.style={draw, rounded corners, draw=black!35, inner sep=3.5mm},
    title/.style={font=\small\bfseries, align=center},
    note/.style={font=\small, text=black!60, align=center},
    arr/.style={-{Latex[length=2mm]}},
]

% --- branded source: Open Food Facts ---
\node[box] (o_bm) at (0mm, 0mm) {Lexical channel\\ (BM25)};
\node[box, right=of o_bm] (o_vec) {Dense channel\\ (embedding + HNSW)};
\coordinate (o_mid) at ($(o_bm)!0.5!(o_vec)$);
\node[title, above=8mm of o_mid] (o_title)
    {Open Food Facts\\ {\normalfont\footnotesize branded products}};
\node[band, below=11mm of o_mid] (o_rrf) {Reciprocal rank fusion};
\node[band, below=of o_rrf] (o_hits) {Top-$k$ branded hits};
\node[frame, fit=(o_title)(o_bm)(o_vec)(o_rrf)(o_hits)] (o_frame) {};

\draw[arr] (o_bm) -- (o_bm |- o_rrf.north);
\draw[arr] (o_vec) -- (o_vec |- o_rrf.north);
\draw[arr] (o_rrf) -- (o_hits);

% --- generic source: USDA FoodData Central ---
\coordinate (u_anchor) at ($(o_frame.east) + (14mm, 0mm)$);
\node[box, anchor=west] (u_bm) at (u_anchor |- o_bm) {Lexical channel\\ (BM25)};
\node[box, right=of u_bm] (u_vec) {Dense channel\\ (embedding + HNSW)};
\coordinate (u_mid) at ($(u_bm)!0.5!(u_vec)$);
\node[title, above=8mm of u_mid] (u_title)
    {USDA FoodData Central\\ {\normalfont\footnotesize generic foods}};
\node[band, below=11mm of u_mid] (u_rrf) {Reciprocal rank fusion};
\node[band, below=of u_rrf] (u_hits) {Top-$k$ generic hits};
\node[frame, fit=(u_title)(u_bm)(u_vec)(u_rrf)(u_hits)] (u_frame) {};

\draw[arr] (u_bm) -- (u_bm |- u_rrf.north);
\draw[arr] (u_vec) -- (u_vec |- u_rrf.north);
\draw[arr] (u_rrf) -- (u_hits);

% --- the one call that fans out into both sources ---
\coordinate (fan) at ($(o_frame.north)!0.5!(u_frame.north)$);
\node[call, anchor=south] (tool) at ($(fan) + (0mm, 16mm)$)
    {Lookup tool call:\\ a batch of queries};
\draw[arr] (tool.west) -| (o_frame.north);
\draw[arr] (tool.east) -| (u_frame.north);

% --- the two ranked lists, returned side by side ---
\coordinate (join) at ($(o_frame.south)!0.5!(u_frame.south)$);
\node[resultbox, anchor=north] (result) at ($(join) + (0mm, -20mm)$)
    {One result block per query,\\ hits grouped by source};
\draw[arr] (o_frame.south) |- (result.west);
\draw[arr] (u_frame.south) |- (result.east);

\node[note, anchor=south, text width=72mm] at ($(result.north) + (0mm, 3mm)$)
    {no fusion across sources};

\end{tikzpicture}

    \caption{Hybrid retrieval: the path of one query in a batch.}
    \label{fig:diagram-retrieval}
\end{figure}

\section{Deterministic totaller}
\label{sec:design-totaller}

When enabled, the totaller will be reachable from two places. The nutrition agent -- and, during the reflection loop, the refiner -- may call it mid-generation to check a proposed plan's running totals before finalising it. Independently of that, the evaluation harness will run the same computation post hoc on every hydrated plan, whether or not the agent itself chose to call the tool. Both call sites will share the same underlying aggregation routine, which will use exact rational arithmetic, so that a total depends on the foods and their grams alone and not on the order the portions happened to be added. The aggregation routine itself will only ever accept a hydrated plan; the wrapper registered as the agent's tool will accept an agent plan and run hydration on it first (\Cref{sec:design-domain-model}), so the number the model sees mid-generation and the number the harness later scores the plan against will, by construction, be computed identically.

\section{Agents and prompt architecture}
\label{sec:design-agents}

The pipeline will use a nutrition agent and, structurally identical to it, a refiner used only inside the reflection loop. When the loop is enabled it will also use a tool-less critic. The critic will list issues; it will not look up foods or revise portions, so a complaint it raises cannot invent a code the later refiner would then have to validate. A single model specification -- a model identifier together with an optional reasoning-effort level, as motivated in \Cref{sec:tools-model-access} -- will be resolved once when a session starts and shared by the nutrition agent, the critic and the refiner. Why the ablation holds critic and generator on the same model, as opposed to giving the reviewer more capability, is taken up in \Cref{sec:impl-critic-design}. The fixed set of dietary guardrails every agent works against will be stated directly in the prompt, because they apply near-universally to a healthy adult; the figures are taken from WHO adult guidance~\cite{who2020healthy,who2012sodium,who2012potassium,whofao2003diet}, with the exception of a per-meal cap on added sugar, stated as a design heuristic that divides the day's allowance across meals so that one dish cannot consume it. Retrieving condition-specific guidance is left as further work (\Cref{sec:conclusions-future-work}). The prompt structure that will carry the profile, the targets and these guardrails to all three agents is the subject of \Cref{sec:design-prompts}.

\section{Prompt design}
\label{sec:design-prompts}

A prompt is ordinary text in the repository and will be governed by the same discipline as the code around it: it will not be duplicated, and it will not be relied upon where the same guarantee can be obtained mechanically. The decisions below apply that principle to the nutrition agent, the critic and the refiner alike.

Each agent's prompt will be split into two tiers. The system prompt will carry everything fixed for the session: the agent's role and an explicit denial of capability, its numbered hard rules, the standing guardrails, the precedence between them, the ordered process, and the output contract. The user turn will carry only the session's data: the free-form request, the profile and the macronutrient targets rendered as JSON, and, for the two reviewing roles, the deterministic totals and the plan under review. The system prompt will thus be a function of the configuration and the user turn a function of the session, so the profile data will never sit among the instructions where a stray field could read as a rule. The role block will state the capability denial in as many words -- the agent has no ability to invent a food or estimate its nutrients -- so that the schema-level restriction of \Cref{sec:design-no-invented-foods} is legible in the prompt and not only in the type it must fill.

The prompts will be zero-shot, carrying no worked example of a finished plan. The output is a large structured object, so an exemplar would cost thousands of tokens on every call and would anchor the model on the specific foods and portions it happened to show -- a contamination the ablation study cannot tolerate, since its variants are meant to differ only by the mechanism toggled. The format-teaching role an exemplar would play will be carried instead by the output schema and its field descriptions.

Nor will any prompt contain an explicit chain-of-thought instruction of the zero-shot kind~\cite{kojima2022large}. Deliberation will be treated as a provider-side control -- the reasoning-effort level of \Cref{sec:tools-model-access}, carried as part of the model identifier -- instead of a phrase in the prompt; what structure the reasoning needs will come from the ordering of the process steps, and the \texttt{rationale} field is an output artefact, not scratch space. Inviting the model to reason in prose about running nutrient sums would invite exactly the free-hand arithmetic the totaller exists to remove.

Where rules conflict, the prompt will state the precedence instead of leaving a flat list for the model to reconcile. A profile condition, a standing guardrail and an explicit user request will not always agree, and sycophancy (\Cref{sec:background-llm-failure-modes}) means an insistent request tends to win regardless of what the prompt says. The prompt will therefore fix when a request legitimately wins -- for the specific item asked for -- and when it does not, requiring the rest of the day to compensate back towards the targets. In the same spirit the guardrails will be phrased as the day's overall tendency, not a per-item ban.

A rule that more than one agent needs will be written once and shared, never restated per agent. The standing guardrail numbers will live in a single fragment that both the nutrition agent's system prompt and the critic's checklist include, and the request-profile-targets head of every user turn will come from one shared composer. Were these duplicated, an edit to one copy -- lowering the sodium ceiling in the planner's prompt but not the critic's, say -- would leave the critic checking against a limit the planner never targeted, surfacing as a fabricated issue in the plan. This is the same discipline that makes the totaller shared code between the runtime and the evaluation harness (\Cref{sec:design-totaller}).

Finally, the prompt will be the last line of defence, not the only one. A contract already enforced mechanically will still be restated in the prompt -- a food code must be copied verbatim from a lookup hit -- not because the sentence guarantees anything, but because stating it lowers how often the output validator has to reject a plan and pay for a retry; the guarantee stays in the validator (\Cref{sec:design-no-invented-foods}). This widens what counts as a prompt: the validator's retry message, the tool descriptions, and the schema's field descriptions are all text the model reads, designed as deliberately as the system prompt itself. The critic's prompt will show the same care -- a reviewer told only to find problems will manufacture them, so an empty issue list will be stated explicitly as the expected outcome for a good plan. The prompt refinements that reading real generated plans forced, together with how the prompt is assembled per ablation variant and for the empty-request case, are taken up in \Cref{sec:impl-prompt-fixes}.

\section{Critique-and-refine loop}
\label{sec:design-reflection}

When enabled, a generated (and validated) plan will be reviewed by the critic--refiner loop already drawn in \Cref{fig:diagram-pipeline}. The critic, given the plan, the user's profile and request, and -- when the totaller is also enabled -- the deterministic totals and any coverage warnings computed for it (\Cref{sec:impl-missing-nutrients}), will report a list of concrete issues, or an empty list if it finds none. Those same totals will also be attached to the refiner's user turn, so both reviewing roles see the same deterministic numbers. An empty list will end the loop immediately with the plan as approved. A nonempty list will be passed to a refiner, structurally identical to the nutrition agent and sharing the same food-lookup and totaller tools, which will be instructed to address only the listed issues with minimal changes elsewhere; its revised plan will then be sent back to the critic for another pass. This will repeat until the critic reports no further issues or a fixed budget of three passes is exhausted; if the budget is exhausted with issues outstanding, the loop will return the last refined plan instead of failing the session outright.

\section{The ablation switch}
\label{sec:design-variant-config}

Because \ref{req:ablatable}~requires every grounding and review mechanism to be independently switchable, the pipeline will expose a single configuration object carrying one boolean flag per switchable mechanism, among them the totaller and the critique-and-refine loop this thesis varies. That object will be read once at the start of a session and threaded through every stage that needs to know whether a given mechanism is active. Food knowledge base lookup is deliberately not one of these flags (\Cref{sec:vision-food-lookup-premise}). This configuration object will be the single source of truth the evaluation harness in \Cref{cha:methodology} sets when it runs the ablation grid, so that a given run's variant label is derived mechanically from exactly the flags that were set for it, not tracked separately and potentially inconsistently. Each agent's system prompt will be assembled from the same flags that decide which tools are registered, so an ablated variant will never describe a tool it no longer has: turning the totaller off will drop the tool and, in the same step, remove the rule that tells the agent to verify its totals.

\section{Package layering}
\label{sec:design-layering}

The codebase will be split into three packages with a strict, one-directional dependency rule. The runtime package, \texttt{dietary\_advisor}, will contain only what is needed to serve a session. The setup package will contain one-off, offline provisioning code -- downloading and building the food knowledge base from the bulk exports chosen in \Cref{sec:tools-food-data}. The evaluation package will contain the ablation harness described in \Cref{cha:methodology}. The setup and evaluation packages may import from the runtime package, but the runtime package may never import from either of them. \Cref{fig:diagram-layering} shows the runtime modules and how the harness reuses them. The two outer packages will lean on the runtime very differently. Setup will reuse the embedding loader, to embed each food record with the same model the session later queries with. The evaluation harness will drive the whole pipeline and reuse hydration and the totaller so that the number it scores a plan against is computed by the same code the session used (\Cref{sec:design-totaller}).

\begin{figure}[htbp]
    \centering
    \begin{tikzpicture}[
    box/.style={draw, rounded corners, align=center, minimum height=7mm, minimum width=34mm, inner sep=1.3mm, font=\small},
    ebox/.style={draw, rounded corners, align=center, minimum height=7mm, minimum width=26mm, inner sep=1.1mm, font=\small},
    arr/.style={-{Latex[length=2mm]}},
]

% --- evaluation (top left) ---
\node[font=\small\bfseries] (etitle) at (0, 0) {evaluation};

\node[ebox] (ecli) at (0, -0.75) {\textbf{cli}};
\node[ebox, minimum width=30mm] (caserunner) at (-3.4, -1.9) {\textbf{case\_runner}};
\node[ebox] (scoring) at (0, -1.9) {\textbf{scoring}};
\node[ebox] (reporting) at (3.4, -1.9) {\textbf{reporting}};
\node[ebox, minimum width=30mm] (validation) at (-2.0, -3.05) {\textbf{validation}};
\node[ebox] (judges) at (2.2, -3.05) {\textbf{judges}};

\begin{scope}[on background layer]
    \node[draw, rounded corners, inner sep=3.2mm,
          fit=(etitle)(ecli)(caserunner)(scoring)(reporting)(validation)(judges)] (evalpkg) {};
\end{scope}

\draw[arr] (ecli) -- (caserunner);
\draw[arr] (ecli) -- (scoring);
\draw[arr] (ecli) -- (reporting);
\draw[arr] (reporting) -- (scoring);
\draw[arr] (scoring) -- (judges);
\draw[arr] (scoring) -- (validation);
\draw[arr] (caserunner) -- (validation);

% --- dietary_advisor (bottom right) ---
\node[font=\small\bfseries] (rtitle) at (5.4, -5.5) {dietary\_advisor};

\node[box] (rcli) at (5.4, -6.25) {\textbf{cli}};
\node[box] (pipeline) at (5.4, -7.35) {\textbf{pipeline}};
\node[box] (reflection) at (5.4, -8.45) {\textbf{reflection}};
\node[box] (agents) at (5.4, -9.55) {\textbf{agents}};
\node[box] (hydration) at (5.4, -10.65) {\textbf{hydration}};

\node[box, minimum width=26mm, minimum height=8.0mm] (profile) at (9.3, -6.25) {\textbf{profile}};
\node[box, minimum width=26mm, minimum height=8.0mm] (telemetry) at (9.3, -9.55) {\textbf{telemetry}};

\node[box, minimum width=26mm] (totaller) at (3.95, -11.95) {\textbf{totaller}};
\node[box, minimum width=26mm] (fooddb) at (6.85, -11.95) {\textbf{food\_db}};

\begin{scope}[on background layer]
    \node[draw, rounded corners, inner sep=3.8mm,
          fit=(rtitle)(rcli)(pipeline)(reflection)(agents)(hydration)(fooddb)(totaller)(profile)(telemetry)] (runtime) {};
\end{scope}

\draw[arr] (rcli) -- (pipeline);
\draw[arr] (pipeline) -- (reflection);
\draw[arr] (reflection) -- (agents);
\draw[arr] (agents) -- (hydration);
\draw[arr] (hydration.south) -- (totaller.north);
\draw[arr] (hydration.south) -- (fooddb.north);
\draw[arr] (pipeline.east) -- ++(0.45, 0) |- ([yshift=1.2mm]agents.east);
\draw[arr] ([yshift=-1.2mm]agents.west) -| ([xshift=-6mm]totaller.north);
\draw[arr] (agents.east) -| ([xshift=6mm]fooddb.north);
\draw[arr] (rcli.east) -- (profile.west);
\draw[arr] (pipeline.east) -| (telemetry.north);

\draw[arr] ([xshift=-7mm]caserunner.south) |- (pipeline.west);
\draw[arr] (validation.south) |- (hydration.west);
\draw[arr] (validation.south) |- (totaller.west);

\end{tikzpicture}

    \caption{Modules of the runtime package and the evaluation package's one-way dependence on them.}
    \label{fig:diagram-layering}
\end{figure}

\section{Reuse of runtime code by the harness}
\label{sec:design-eval-reuse}

The harness's trustworthiness will rest on the division between the runtime code it reuses and the evaluation logic it owns. Its case runner will drive the very same pipeline a real session uses, and its validation stage will recompute each plan's totals with the same hydration and totaller code, so the harness will never re-implement any logic whose result it is trying to measure. Everything the runtime must not know about -- the scenarios, the judges, the metrics, and the reporting -- the harness will own itself. The judge, by contrast, will be handed a plan carrying no nutrient numbers and no indication of which variant or model produced it, scored against targets frozen on the profile outside the pipeline under test.
